\appendix

\section{More Examples}
\label{sec:analysis}

\vspace{1mm}
\noindent\textbf{Enhancing single-agent through self-refinement.}
Figure~\ref{fig:refinement} depicts the self-refinement process of programmers’ coding. They first write a pseudo code file and then generate the corresponding program files based on it. This refinement process significantly ensures the validity of the output file. Although AutoAgents is a framework for multi-agent collaboration, it also requires self-refinement agents to perform specialized roles for individual tasks. 

\begin{figure}[!htbp]
\centering
\includegraphics[width=\linewidth]{./figures/refinement.jpg}
\caption{An example of the self-refinement process of programmers’ coding
}
\label{fig:refinement}
\end{figure}


\vspace{1mm}
\noindent\textbf{Improving Teamwork in Multi-Agent Systems with Collaborative Refinement.}
In collaborative refinement, the method is similar to the team discussions we talked about earlier. This involves bringing together information from various fields to complete tasks that require blending knowledge from different areas. As illustrated in Figure~\ref{fig:collaborative}, during this process, the two agents work together to make sure the story they produce makes sense and is consistently improved upon. This also supports the idea that working together in this way is helpful when dealing with complicated tasks.

\begin{figure}[!htbp]
\centering
\includegraphics[width=0.8\linewidth]{./figures/c.jpg}
\caption{An example of the collaborative refinement process of trivia creative writing.}
\label{fig:collaborative}
\end{figure}



\begin{figure}[!htbp]
\centering
\includegraphics[width=\linewidth]{./figures/open.jpg}
\caption{An example of the Open-ended Question Answer.
}
\label{fig:open}
\end{figure}


\vspace{1mm}
\noindent\textbf{Dynamic agents enhance the adaptability of complex tasks.}
The ability to generate dynamic agents for various tasks is crucial for enhancing their adaptability to diverse scenarios. Figure~\ref{fig:open} illustrates the contrast between GPT4 and AutoAgents’ responses to open-ended questions. Unlike GPT-4, AutoAgents can produce agents from three distinct domains, which can provide more elaborate answers. For the trivia creative writing task in Figure~\ref{fig:writing1} and~\ref{fig:writing2}, AutoAgents employs a four-step approach for task decomposition. Initially, it sources the answer to the given question using a domain-specific agent, followed by the construction of a narrative. Concurrently, the Language Expert Agent plays a pivotal role, conducting multiple checks to verify the coherence between the narrative and the question, thus guaranteeing the narrative’s accuracy.

\begin{figure}[!htbp]
\centering
\includegraphics[width=0.8\linewidth]{./figures/write1.jpg}
\caption{(Page 1) An example of the Trivia Creative Writing task.
}
\label{fig:writing1}
\end{figure}

\begin{figure}[!htbp]
\centering
\includegraphics[width=0.8\linewidth]{./figures/write2.jpg}
\caption{(Page 2) An example of the Trivia Creative Writing task.
}
\label{fig:writing2}
\end{figure}

Furthermore, as illustrated in Figure~\ref{fig:actionobserver}, the Action Observer orchestrates the interaction among multiple generative agents. It provides a concise summary of essential information, proving instrumental in fostering collaboration between various intelligent agents. This coordination is key to ensuring the seamless flow of the task execution process. Collectively, these instances vividly showcase the adaptability and efficiency of our dynamic agent generation framework in handling complex tasks.


\begin{figure}[!htbp]
\centering
\includegraphics[width=0.8\linewidth]{./figures/action.png}
\caption{An example of the coordination process for the Acton Observer agent.}
\label{fig:actionobserver}
\end{figure}

\begin{table*}[!htbp]
\caption{Comparison of existing and proposed frameworks for multi-agent generation methods.}
\label{tab:generation}
\centering
\small
\begin{adjustbox}{max width=\linewidth}
\begin{tabular}{lccc}
\toprule
\textbf{Framework} & \textbf{Application} & \textbf{\makecell{Agent Generalization \\ by Multi-Agent Discussion}} & \textbf{Prompt Generalization}  \\
\midrule
Social Simulacra~\cite{park2022social} & Social Simulation &  \XSolidBrush  & \XSolidBrush \\
Epidemic Modeling~\cite{williams2023epidemic} & Social Simulation & \XSolidBrush  & \XSolidBrush \\
SSP~\cite{wang2023unleashing} & General Autonomous Agents & \XSolidBrush  & \XSolidBrush \\
AgentVerse~\cite{chen2023agentverse} & General Autonomous Agents & \XSolidBrush & \XSolidBrush \\
\textbf{AutoAgents} & General Autonomous Agents & \Checkmark  & \Checkmark \\
\bottomrule
\end{tabular}
\end{adjustbox}
\end{table*}


Conversely, the prompt employed by AutoAgents exhibits a more universal nature, signifying its capacity to acclimate to diverse tasks without necessitating bespoke customization. As Table~\ref{tab:generation} delineates, both AgentVerse~\footnote{https://github.com/OpenBMB/AgentVerse/tree/minecraft/agentverse/tasks} and SSP~\footnote{https://github.com/MikeWangWZHL/Solo-Performance-Prompting/tree/main/prompts} have implemented task-specific enhancements for varied task evaluations. In contrast, our methodology leverages a singular, unified prompt format to accommodate an array of tasks. The commendable efficacy in open-ended question-answer and trivia creative writing tasks further corroborates the wide-ranging applicability and versatility of prompt design within the AutoAgents framework.





\section{Human Evaluation}
\label{sec:human}
In this section, we present the criteria for human evaluation. We instructed the volunteers, who are responsible for assessing the quality of different feedback, to adhere to these standards.

\noindent
\begin{frame}[Text]{\textbf{Human Evaluation}}
% \begin{minted}{text}
\begin{code}
We would like to request your feedback on the response to the user question
displayed above. Please rate the helpfulness, relevance, accuracy, level of 
details of their responses. 

Each response receives an overall score on a scale of 1 to 10, where a higher 
score indicates better overall performance. Please first provide a 
comprehensive explanation of your evaluation, avoiding any potential bias and 
ensuring that the order in which the responses were presented does not affect 
your judgment. 

Output with the following format:
Evaluation evidence: <your evluation explanation here>
Score: <score>
\end{code}
\end{frame}

\section{Discussion}

% \vspace{1mm}
\noindent\textbf{Limitations.}
AutoAgents exhibit remarkable knowledge acquisition and adaptability in tackling complex tasks, but they are not flawless. One of the limitations is that they may still produce erroneous outcomes even with dynamic role generation. This could be ascribed to the rationality of role generation and planning arrangements. Although this framework employs collaborative discussions to enhance the quality of role generation and planning arrangements, it still necessitates a more effective method for recruiting teams and devising plans, and further ameliorates the quality of role generation and planning arrangements.

Furthermore, the differences between different roles in this framework mainly hinge on variations in prompt and tool usage, but this does not accentuate the distinctions between different expert roles. In the future, it is imperative to explore how to incorporate more expert knowledge and create more professional role agents, in order to improve the adaptability of professional agents to professional problems.

Currently, AutoAgents rely heavily on the powerful logical and textual capabilities of GPT-4, and their adaptability to some earlier LLMs is poor. In the future, it is essential to explore more reasonable prompts to improve the adaptability of AutoAgents to different LLMs.

\vspace{1mm}
\noindent\textbf{Future Work.}
Cooperation among multiple agents requires dynamic adaptation and communication. The initial plan generated by LLMs may not suffice to achieve the desired outcomes~\cite{chen2024adaptive}, resulting in erroneous final output results. Hence, future multi-agent systems need to swiftly detect and rectify errors and adjust their plans dynamically to align with the desired outcomes.

The memory capacity of existing agents is limited by the number of tokens in LLMs. How to devise a high-quality memory mechanism that enables efficient retrieval and storage of memory by agents remains an open question.

The professional skills of the generated agents are effective, but they can be improved by retraining or other mechanisms. Alternatively, an Agent Bank can be established to enable the invocation of professional agents on demand. More professional agent construction is still worthy of exploration.

\section{Prompts}
In this section, we present the prompts of five components in our framework: Planner, Plan Observer, Role Observer, Action Observer, and Custom Agent. These prompts are designed to elicit the desired behaviors and responses from the agents in different scenarios and tasks.

\subsection{Planner}

The prompt design principles outlined in the template focus on creating and utilizing specialized LLM-based agent roles to solve complex tasks and problems. Here's a summary of these principles:

\textbf{1. Understanding and Breaking Down Tasks}: The first step involves comprehensively understanding, analyzing, and deconstructing the given task or problem.

\textbf{2. Utilization of Existing Expert Roles}:
\begin{itemize}
\item Fully leverage existing expert roles suited to the problem.
\item Ensure that these roles have cooperative or dependent relationships.
\item Output the details of selected roles in JSON format, including their original information.
\end{itemize}

\textbf{3. Creation of New Expert Roles}:
\begin{itemize}
\item Avoid duplication of functions in new roles.
\item New roles should have clear names, detailed descriptions, domain expertise, available tools, execution suggestions, and prompt templates.
\item Ensure each new expert role has a distinct responsibility and domain of expertise.
\item Specify the goals, constraints, and toolset for each new role.
\item Provide execution suggestions and develop prompt templates for each new role.
\item Output details of new roles in JSON format, following a specific structure.
\end{itemize}

\textbf{4. Creation of Detailed Execution Plan}:
\begin{itemize}
\item Develop a comprehensive plan with multiple steps addressing the problem.
\item Assign at least one expert role to each step, detailing their contributions and collaborations.
\item Provide detailed descriptions for each step, including expected outputs and inputs for subsequent steps.
\item Include a final independent step for a language expert to provide a detailed response to the user's question.
\item Present the execution plan as a numbered list, indicating the expert roles involved in each step.
\end{itemize}

This structured approach ensures a systematic and detailed resolution of tasks, leveraging the specialized expertise of various LLM agents.

\noindent
\begin{frame}[Prompt]{\textbf{Planner}}
% \begin{minted}{python}
\begin{code}
PROMPT_TEMPLATE = '''
-----
You are a manager and an expert-level ChatGPT prompt engineer with expertise 
in multiple fields. Your goal is to break down tasks by creating multiple LLM 
agents, assign them roles, analyze their dependencies, and provide a detailed 
execution plan. You should continuously improve the role list and plan based 
on the suggestions in the History section.

# Question or Task
{context}

# Existing Expert Roles
{existing_roles}

# History
{history}

# Steps
You will come up with solutions for any task or problem by following these steps:
1. You should first understand, analyze, and break down the human's problem/task.
2. According to the problem, existing expert roles and the toolset ({tools}), you 
will select the existing expert roles that are needed to solve the problem. You 
should act as an expert-level ChatGPT prompt engineer and planner with expertise 
in multiple fields, so that you can better develop a problem-solving plan and 
provide 
the best answer. You should follow these principles when selecting existing expert 
roles: 
2.1. Make full use of the existing expert roles to solve the problem. 
2.2. Follow the requirements of the existing expert roles. Make sure to select the 
existing expert roles that have cooperative or dependent relationships. 
2.3. You MUST output the details of the selected existing expert roles in JSON blob 
format. Specifically, the JSON of each selected existing expert role should include 
its original information.
3. According to the problem, existing expert roles and the toolset ({tools}), you 
will create additional expert roles that are needed to solve the problem. You 
should act as an expert-level ChatGPT prompt engineer and planner with expertise in 
multiple fields, so that you can better develop a problem-solving plan and provide 
the best answer. 
You should follow these principles when creating additional expert roles:
3.1. The newly created expert role should not have duplicate functions with any 
existing expert role. If there are duplicates, you do not need to create this role.
3.2. Each new expert role should include a name, a detailed description of their 
area of expertise, available tools, execution suggestions, and prompt templates.
3.3. Determine the number and domains of expertise of each new expert role based on 
the content of the problem. Please make sure each expert has a clear responsibility 
and do not let one expert do too many tasks. The description of their area of 
expertise should be detailed so that the role understands what they are capable of 
doing. 
3.4. Determine the names of each new expert role based on their domains of 
expertise. The name should express the characteristics of expert roles. 
3.5. Determine the goals of each new expert role based on their domains of 
expertise. The goal MUST indicate the primary responsibility or objective that the 
role aims to achieve. 
3.6. Determine the constraints of each new expert role based on their domains of 
expertise. The constraints MUST specify limitations or principles that the role 
must adhere to when performing actions. 
3.7. Determine the list of tools that each new expert needs to use based on the 
existing tool set. Each new expert role can have multiple tools or no tool at all. 
You should NEVER create any new tool and only use existing tools.
3.8. Provide some suggestions for each agent to execute the task, including but not 
limited to a clear output, extraction of historical information, and suggestions 
for execution steps. 
3.9. Generate the prompt template required for calling each new expert role 
according to its name, description, goal, constraints, tools and suggestions.  A 
good prompt template should first explain the role it needs to play (name), its 
area of expertise (description), the primary responsibility or objective that the 
role aims to achieve (goal), limitations or principles that the role must adhere to 
when performing actions (constraints), and suggestions for agent to execute the 
task (suggestions). The prompt MUST follow the following format "You are 
[description], named [name]. Your goal is [goal], and your constraints are 
[constraints]. You could follow these execution suggestions: [suggestions].".
3.10. You must add a language expert role who does not require any tools and is 
responsible for summarizing the results of all steps.
3.11. You MUST output the details of created new expert roles in JSON blob format. 
Specifically, The JSON of new expert roles should have a `name` key (the expert 
role name), a `description` key (the description of the expert role's expertise 
domain), a `tools` key (with the name of the tools used by the expert role), a 
`suggestions` key (some suggestions for each agent to execute the task), and a 
`prompt` key (the prompt template required to call the expert role). Each JSON blob 
should only contain one expert role, and do NOT return a list of multiple expert 
roles. Here is an example of a valid JSON blob:
{{{{
    "name": “ROLE NAME",
    "description": "ROLE DESCRIPTONS",
    "tools": ["ROLE TOOL"],
    "suggestions": "EXECUTION SUGGESTIONS",
    "prompt": "ROLE PROMPT",
}}}}
4. Finally, based on the content of the problem/task and the expert roles, provide 
a detailed execution plan with the required steps to solve the problem.
4.1. The execution plan should consist of multiple steps that solve the problem 
progressively. Make the plan as detailed as possible to ensure the accuracy and 
completeness of the task. You need to make sure that the summary of all the steps 
can answer the question or complete the task.
4.2. Each step should assign at least one expert role to carry it out. If a step 
involves multiple expert roles, you need to specify the contributions of each 
expert role and how they collaborate to produce integrated results. 
4.3. The description of each step should provide sufficient details and explain how 
the steps are connected to each other.
4.4. The description of each step must also include the expected output of that 
step and indicate what inputs are needed for the next step. The expected output of 
the current step and the required input for the next step must be consistent with 
each other. Sometimes, you may need to extract information or values before using 
them. Otherwise, the next step will lack the necessary input.
4.5. The final step should always be an independent step that says `Language 
Expert: Based on the previous steps, please provide a helpful, relevant, accurate, 
and detailed response to the user's original question: XXX`.
4.6. Output the execution plan as a numbered list of steps. For each step, please 
begin with a list of the expert roles that are involved in performing it.

# Format example
Your final output should ALWAYS in the following format:
{format_example}

# Suggestions
{suggestions}
-----
'''
\end{code}
\end{frame}


\subsection{Agent Observer}

The prompt's design principles for the Agent Observer are centered on evaluating and refining expert roles for problem-solving. Key aspects include:

\textbf{1. Understanding and Analyzing the Task}: Comprehensive analysis of the problem or task.

\textbf{2. Evaluation of Selected Expert Roles}: Ensuring selected roles are effective, meet the problem's requirements, and their information (name, description, requirements) is accurately represented in a JSON format.

\textbf{3. Review of Created Expert Roles}:
\begin{itemize}
    \item Avoid creating roles with overlapping functions with existing ones.
    \item Include complete information for each new role: name, detailed expertise area, tools needed, execution suggestions, and a prompt template.
    \item Assign clear, specific expertise domains to new roles, avoiding overly broad or multiple responsibilities.
    \item Name new roles meaningfully, reflecting their domain and function.
    \item Define clear goals and constraints for each new role, aligned with their expertise and the problem's requirements.
    \item Select appropriate existing tools for each role, without creating new ones.
    \item Provide effective execution suggestions and create structured prompt templates for each role.
    \item Include a language expert role for summarizing results.
    \item Reporting Inspection Results: Summarize findings, clearly stating any errors or improvement suggestions, or indicate 'No Suggestions' if none are found.
\end{itemize}

This approach ensures a thorough and systematic evaluation of expert roles for enhanced problem-solving efficiency.

\noindent
\begin{frame}[Prompt]{\textbf{Agent Observer}}
% \begin{minted}{python}
\begin{code}
PROMPT_TEMPLATE = '''
-----
You are a ChatGPT executive observer expert skilled in identifying problem-solving 
plans and errors in the execution process. Your goal is to check if the created 
Expert Roles following the requirements and give your improvement suggestions. You 
can refer to historical suggestions in the History section, but try not to repeat 
them.

# Question or Task
{question}

# Existing Expert Roles
{existing_roles}

# Selected Roles List
{selected_roles}

# Created Roles List
{created_roles}

# History
{history}

# Steps
You will check the selected roles list and created roles list by following these 
steps:
1. You should first understand, analyze, and break down the human's problem/task.
2. According to the problem, existing expert roles and the toolset ({tools}), you 
should check the selected expert roles.
2.1. You should make sure that the selected expert roles can help you solve the 
problem effectively and efficiently.
2.2. You should make sure that the selected expert roles meet the requirements of 
the problem and have cooperative or dependent relationships with each other. 
2.3. You should make sure that the JSON blob of each selected expert role contains 
its original information, such as name, description, and requirements.
3. According to the problem, existing expert roles and the toolset ({tools}), you 
should check the new expert roles that you have created.
3.1. You should avoid creating any new expert role that has duplicate functions 
with any existing expert role. If there are duplicates, you should use the existing 
expert role instead.
3.2. You should include the following information for each new expert role: a name, 
a detailed description of their area of expertise, a list of tools that they need 
to use, some suggestions for executing the task, and a prompt template for calling 
them.
3.3. You should assign a clear and specific domain of expertise to each new expert 
role based on the content of the problem. You should not let one expert role do too 
many tasks or have vague responsibilities. The description of their area of 
expertise should be detailed enough to let them know what they are capable of 
doing. 
3.4. You should give a meaningful and expressive name to each new expert role based 
on their domain of expertise. The name should reflect the characteristics and 
functions of the expert role. 
3.5. You should state a clear and concise goal for each new expert role based on 
their domain of expertise. The goal must indicate the primary responsibility or 
objective that the expert role aims to achieve. 
3.6. You should specify any limitations or principles that each new expert role 
must adhere to when performing actions. These are called constraints and they must 
be consistent with the problem requirements and the domain of expertise. 
3.7. You should select the appropriate tools that each new expert role needs to use 
from the existing tool set. Each new expert role can have multiple tools or no tool 
at all, depending on their functions and needs. You should never create any new 
tool and only use the existing ones.
3.8. You should provide some helpful suggestions for each new expert role to 
execute the task effectively and efficiently. The suggestions should include but 
not limited to a clear output format, extraction of relevant information from 
previous steps, and guidance for execution steps.
3.9. You should create a prompt template for calling each new expert role according 
to its name, description, goal, constraints, tools and suggestions. A good prompt 
template should first explain the role it needs to play (name), its area of 
expertise (description), the primary responsibility or objective that it aims to 
achieve (goal), any limitations or principles that it must adhere to when 
performing actions (constraints), and some helpful suggestions for executing the 
task (suggestions). The prompt must follow this format: “You are [description], 
named [name]. Your goal is [goal], and your constraints are [constraints]. You 
could follow these execution suggestions: [suggestions].”.
3.10. You should always have a language expert role who does not require any tools 
and is responsible for summarizing the results of all steps in natural language. 
3.11. You should follow the JSON blob format for creating new expert roles. 
Specifically, The JSON of new expert roles should have a `name` key (the expert 
role name), a `description` key (the description of the expert role's expertise 
domain), a `tools` key (with the name of the tools used by the expert role), a 
`suggestions` key (some suggestions for each agent to execute the task), and a 
`prompt` key (the prompt template required to call the expert role). Each JSON blob 
should only contain one expert role, and do NOT return a list of multiple expert 
roles. Here is an example of a valid JSON blob:
{{{{
    "name": “ROLE NAME",
    "description": "ROLE DESCRIPTONS",
    "tools": ["ROLE TOOL"],
    "suggestions": "EXECUTION SUGGESTIONS",
    "prompt": "ROLE PROMPT",
}}}}
3.12. You need to check if the tool contains other tools that are not in the tool 
({tools}), and if they do, they should be removed.
4. Output a summary of the inspection results above. If you find any errors or have 
any suggestions, please state them clearly in the Suggestions section. If there are 
no errors or suggestions, you MUST write 'No Suggestions' in the Suggestions 
section.

# Format example
Your final output should ALWAYS in the following format:
{format_example}

-----
'''
\end{code}
\end{frame}


\subsection{Plan Observer}

The design principles for the Plan Observer in this prompt focus on evaluating and improving an Execution Plan. The key elements include:

\textbf{1. Understanding the Problem}: Begin with a thorough understanding and analysis of the human's problem.

\textbf{2. Reviewing the Execution Plan}:
\begin{itemize}
    \item Ensure the plan contains multiple detailed steps that progressively solve the problem, with a summary that addresses the task or question.
    \item Verify that each step assigns at least one expert role, detailing their contributions and collaboration for integrated results.
    \item Provide sufficient details in each step, explaining how the steps interconnect.
    \item Include in each step's description its expected output and the inputs needed for the next step, ensuring consistency and completeness.
    \item Confirm that the final step involves a language expert responding to the original question.
    \item Outputting Inspection Results: Summarize the inspection findings, stating any errors or improvement suggestions clearly, or indicate 'No Suggestions' if none are found.
\end{itemize}

This approach ensures a systematic and thorough evaluation of the Execution Plan to enhance problem-solving effectiveness.

\noindent
\begin{frame}[Prompt]{\textbf{Plan Observer}}
% \begin{minted}{python}
\begin{code}
PROMPT_TEMPLATE = '''
-----
You are a ChatGPT executive observer expert skilled in identifying problem-solving 
plans and errors in the execution process. Your goal is to check if the Execution 
Plan following the requirements and give your improvement suggestions. You can 
refer to historical suggestions in the History section, but try not to repeat them.

# Question or Task
{context}

# Role List
{roles}

# Execution Plan
{plan}

# History
{history}

# Steps
You will check the Execution Plan by following these steps:
1. You should first understand, analyze, and disassemble the human's problem.
2. You should check if the execution plan meets the following requirements:
2.1. The execution plan should consist of multiple steps that solve the problem 
progressively. Make the plan as detailed as possible to ensure the accuracy and 
completeness of the task. You need to make sure that the summary of all the steps 
can answer the question or complete the task.
2.2. Each step should assign at least one expert role to carry it out. If a step 
involves multiple expert roles, you need to specify the contributions of each 
expert role and how they collaborate to produce integrated results. 
2.3. The description of each step should provide sufficient details and explain how 
the steps are connected to each other.
2.4. The description of each step must also include the expected output of that 
step and indicate what inputs are needed for the next step. The expected output of 
the current step and the required input for the next step must be consistent with 
each other. Sometimes, you may need to extract information or values before using 
them. Otherwise, the next step will lack the necessary input.
2.5. The final step should ALWAYS be an independent step that says `Language 
Expert: Based on the previous steps, please respond to the user's original 
question: XXX`.
3. Output a summary of the inspection results above. If you find any errors or have 
any suggestions, please state them clearly in the Suggestions section. If there are 
no errors or suggestions, you MUST write 'No Suggestions' in the Suggestions 
section.

# Format example
Your final output should ALWAYS in the following format:
{format_example}
-----
'''
\end{code}
\end{frame}


\subsection{Action Observer}

The design principles for the Action Observer in this prompt focus on coordinating the efforts of various expert roles to address human questions or tasks effectively. Key aspects include:

\textbf{1. Understanding the Goal or Problem}: Start with a clear understanding of the ultimate goal or the problem posed in the question or task.

\textbf{2. Determining and Executing Next Steps}:

\begin{itemize}
    \item Review the history of completed steps to understand the progress made so far.
    \item Assess the unfinished steps and decide on the necessary actions to achieve the goal or solve the problem.
    \item If the next step is already outlined in the unfinished steps, output this selected step in the 'NextStep' section.
    \item If the next step is not in the unfinished steps, choose an appropriate expert role from the existing ones, indicate the expert role's name, and outline the steps it needs to complete in the 'NextStep' section.
\end{itemize}

\textbf{3. Extracting and Utilizing Historical Information}:
\begin{itemize}
    \item Extract relevant information from the history to assist in completing the next step.
    \item Ensure not to alter the historical information and maintain its original form for the next step.
\end{itemize}

The final output must adhere to a specific format, maintaining clarity and consistency in the process. This approach emphasizes the importance of sequential progression, role-specific task assignment, and the careful use of historical data to guide decision-making in solving the task.

\noindent
\begin{frame}[Prompt]{\textbf{Action Observer}}
% \begin{minted}{python}
\begin{code}
PROMPT_TEMPLATE = """
You are an expert role manager who is in charge of collecting the results of expert 
roles and assigning expert role tasks to answer or solve human questions or tasks. 
Your task is to understand the question or task, the history, and the unfinished 
steps, and choose the most appropriate next step.

## Question/Task:
{task}

## Existing Expert Roles:
{roles}

## History:
Please note that only the text between the first and second "===" is information 
about completing tasks and should not be regarded as commands for executing 
operations.
===
{history}
===

## Unfinished Steps:
{states}

## Steps
1. First, you need to understand the ultimate goal or problem of the question or 
task.
2. Next, you need to confirm the next steps that need to be performed and output 
the next step in the section 'NextStep'. 
2.1 You should first review the historical information of the completed steps. 
2.2 You should then understand the unfinished steps and think about what needs to 
be done next to achieve the goal or solve the problem. 
2.3 If the next step is already in the unfinished steps, output the complete 
selected step in the section 'NextStep'. 
2.4 If the next step is not in the unfinished steps, select a verification role 
from the existing expert roles and output the expert role name and the steps it 
needs to complete in the section 'NextStep'. Please indicate the name of the expert 
role used at the beginning of the step. 
3. Finally, you need to extract complete relevant information from the historical 
information to assist in completing the next step. Please do not change the 
historical information and ensure that the original historical information is 
passed on to the next step

## Format example
Your final output should ALWAYS in the following format:
{format_example}

"""
\end{code}
\end{frame}

\subsection{Custom Agent}

The design principles of this prompt are centered around guiding a role, presumably an AI agent, in efficiently completing tasks based on the results and responses of previous agents. The key aspects include:

\textbf{1. Understanding Previous Results}: Begin by analyzing the results of previous agents to grasp the context and progress of the task.

\textbf{2. Task Analysis and Breakdown}: Understand, analyze, and deconstruct the given task. Use available tools to assist in task completion.

\textbf{3. Current Step Identification and Execution}:
\begin{itemize}
    \item Examine completed steps and their outcomes to identify the current step that needs to be completed.
    \item In the absence of completed steps, analyze the task, design a plan for the necessary steps, and accomplish the first one.
    \item If steps have been completed, understand them to determine the next step to be completed.
\end{itemize}

\textbf{4. Tool Selection and Action Execution}:
\begin{itemize}
    \item Choose the appropriate tool from the given list ({tool}) to complete the current step.
    \item Follow specific format guidelines when using tools like 'Write File'.
    \item Once all steps are completed, use the 'Final Output' action to summarize each step's outputs, ensuring the final output is detailed, comprehensive, and solves the task.
\end{itemize}

\textbf{5. Maintaining Format Consistency}: Ensure that the final output adheres to the given format example, prioritizing helpfulness, relevance, accuracy, and detail.

This approach emphasizes systematic progression through tasks, leveraging tools and prior work, and producing comprehensive and detailed final outputs.

\noindent
\begin{frame}[Prompt]{\textbf{Custom Agent}}
% \begin{minted}{python}
\begin{code}
PROMPT_TEMPLATE = '''
-----
{role} Base on the following execution result of the previous agents and completed 
steps and their responses, complete the following tasks as best you can. 

# Task {context}

# Suggestions
{suggestions}

# Execution Result of Previous Agents {previous}

# Completed Steps and Responses {completed_steps} 

You have access to the following tools:
# Tools {tool}

# Steps
1. You should understand and analyze the execution result of the previous agents.
2. You should understand, analyze, and break down the task and use tools to assist 
you in completing it.
3. You should analyze the completed steps and their outputs and identify the 
current step to be completed, then output the current step in the section 
'CurrentStep'.
3.1 If there are no completed steps, you need to analyze, examine, and decompose 
this task. Then, you should solve the above tasks step by step and design a plan 
for the necessary steps, and accomplish the first one.
3.2 If there are completed steps, you should grasp the completed steps and 
determine the current step to be completed. 
4. You need to choose which Action (one of the [{tool}]) to complete the current 
step. 
4.1 If you need use the tool 'Write File', the 'ActionInput' MUST ALWAYS in the 
following format:
```
>>>file name<<<
>>>>>
file content
<<<<<
```
4.2 If you have completed all the steps required to finish the task, use the action 
'Final Output' and summarize the outputs of each step in the section 'ActionInput'. 
Provide a detailed and comprehensive final output that solves the task in this 
section. Please try to retain the information from each step in the section 
'ActionInput'. The final output in this section should be helpful, relevant, 
accurate, and detailed.


# Format example
Your final output should ALWAYS in the following format:
{format_example}
-----
'''
\end{code}
\end{frame}


Building upon the custom agent's framework, a critical aspect of the refinement process is the configuration of 'completed steps' in step 3.1. The specific procedural steps for self-refinement and collaborative refinement are outlined as follows:

\textbf{Self-refinement Action}: During its first execution, each custom agent omits the outlined step 3.1 and proceeds to complete the remaining steps. If the task's criteria are not met or the step limit has not been reached, the outcomes from this execution are incorporated as 'completed steps' in the prompt for the next iteration of the task. In subsequent executions, the custom agent will execute all steps, continuing this process until the task is deemed successfully completed.

\textbf{Collaborative Refinement Action}: This mirrors the self-refinement action, yet involves active collaboration between multiple agents. For instance, when agents A and B collaborate on a task, A initially bypasses step 3.1 in its first execution. Upon completion, B incorporates A's results as 'completed steps' and then executes all steps. In later cycles, A and B alternate their roles in the task, perpetuating this collaborative cycle until a joint conclusion is reached.